\section{Self-Organizing Map Methodology}
\label{sec:som_methodology}

\subsection{Purpose of the SOM}

A Self-Organizing Map (SOM) was introduced to investigate whether
unsupervised learning could identify physically meaningful flow states in the
CFD solution and whether information derived from those states could
subsequently contribute to mesh-refinement decisions.

The SOM is an unsupervised neural-network method that maps multidimensional
input data onto a lower-dimensional discrete lattice while attempting to
preserve topological relationships in the input space
\cite{kohonen1982,kohonen2001}. SOM-based clustering and interpretation of
the trained map have subsequently been developed for exploratory analysis of
multidimensional data \cite{vesanto2000}.

In the present work, the SOM is not used as a surrogate for the CFD solver.
The Navier--Stokes equations remain the source of the physical solution.
Instead, the SOM analyzes quantities extracted from the CFD field.

The initial SOM analysis is performed using the M40 solution, containing
1600 cells. Importantly, the M160 reference solution is not used during SOM
training or during the construction of the refinement indicators. It is
reserved for subsequent evaluation.

The supporting data-analysis workflow was implemented in Python using the
scientific-computing ecosystem. Array and numerical operations relied on
NumPy and SciPy, while preprocessing utilities were drawn from
scikit-learn \cite{harris2020,virtanen2020,pedregosa2011}.

\subsection{Candidate physical features}

For each M40 cell, the CFD post-processing pipeline extracts the cell-center
coordinates and local flow quantities. The candidate nondimensional
quantities include

\begin{equation}
\frac{U_x}{U_{\mathrm{lid}}},
\qquad
\frac{U_y}{U_{\mathrm{lid}}},
\qquad
\frac{|\mathbf{U}|}{U_{\mathrm{lid}}},
\qquad
\frac{L|\nabla\mathbf{U}|}{U_{\mathrm{lid}}},
\qquad
\frac{L|\boldsymbol{\omega}|}{U_{\mathrm{lid}}},
\end{equation}

where

\begin{equation}
\boldsymbol{\omega}=\nabla\times\mathbf{U}.
\end{equation}

The velocity-gradient magnitude is computed using the Frobenius norm of the
velocity-gradient tensor.

Spatial coordinates are deliberately excluded from the SOM input vector.
Consequently, any spatially coherent classification produced by the SOM must
arise from similarities in the local flow state rather than from explicit
knowledge of cell position.

\subsection{Feature selection}

An exploratory correlation analysis showed a strong correlation between the
velocity-gradient magnitude and the vorticity magnitude,

\begin{equation}
r\left(|\nabla\mathbf{U}|,|\boldsymbol{\omega}|\right)
=
0.981.
\end{equation}

To reduce this strong redundancy, vorticity was excluded from the first SOM
input vector.

The resulting feature vector was

\begin{equation}
\mathbf{x}_i =
\begin{bmatrix}
U_{x,i}^{*} &
U_{y,i}^{*} &
|\mathbf{U}_i|^{*} &
|\nabla\mathbf{U}_i|^{*}
\end{bmatrix}^{T},
\end{equation}

where the superscript \(^{*}\) denotes physical nondimensionalization using
the characteristic cavity length and lid velocity.

The velocity magnitude is mathematically related to its two velocity
components and therefore introduces some redundancy. It was retained in this
initial investigation to preserve an explicit scalar measure of local flow
speed. Feature-ablation studies are left for future work.

\subsection{Standardization}

Physical nondimensionalization and machine-learning standardization are
treated as separate operations.

After nondimensionalization, each SOM feature is standardized according to

\begin{equation}
z_j =
\frac{x_j-\mu_j}{\sigma_j},
\end{equation}

where \(\mu_j\) and \(\sigma_j\) are respectively the mean and standard
deviation of feature \(j\) over the M40 dataset.

This prevents variables with larger numerical ranges from dominating the
Euclidean distance used during SOM training.

\subsection{SOM architecture and training}

The SOM consists of a \(4\times4\) rectangular lattice containing 16 neurons.
The implementation was developed in Python using NumPy, with feature
standardization performed using \texttt{StandardScaler}.

Training used 20,000 iterations with a fixed random seed of 42 to provide
reproducibility. The learning rate decreased exponentially from

\begin{equation}
\alpha_0=0.5
\end{equation}

to

\begin{equation}
\alpha_f=0.05,
\end{equation}

while the neighborhood radius decreased from

\begin{equation}
\sigma_0=2.0
\end{equation}

to

\begin{equation}
\sigma_f=0.5.
\end{equation}

A Gaussian neighborhood function was used during weight updates.

For each CFD cell, the best-matching unit (BMU) is the neuron whose weight
vector has the smallest Euclidean distance to the standardized feature
vector.

\subsection{Quantization error}

The cell-level SOM quantization error is defined as

\begin{equation}
QE_i =
\left\|
\mathbf{z}_i-\mathbf{w}_{b(i)}
\right\|_2,
\end{equation}

where \(\mathbf{z}_i\) is the standardized feature vector of cell \(i\) and
\(\mathbf{w}_{b(i)}\) is the weight vector of its BMU.

For the trained M40 SOM, the mean cell-level quantization error was

\begin{equation}
\overline{QE}=0.4976,
\end{equation}

while the maximum value was

\begin{equation}
QE_{\max}=7.5142.
\end{equation}

All 16 neurons contained assigned CFD cells after training.

A large quantization error indicates that a CFD state is represented
relatively poorly by its assigned SOM neuron. It does not, by definition,
represent CFD discretization error. Distinguishing these two quantities is
central to the experiments presented later in this report.

\subsection{Interpretation of SOM states}

Despite the exclusion of the spatial coordinates from the training features,
the resulting BMU assignments form spatially coherent regions when mapped
back onto the cavity.

This observation indicates that the SOM is identifying recurring local flow
states rather than arbitrary spatial locations. Different neurons represent
combinations of velocity direction, velocity magnitude, and velocity-gradient
magnitude.

The SOM therefore provides a classification of the local physical state.
Whether this classification contains information useful for adaptive
refinement is treated as an experimental question rather than an assumption.
